\properlane{gift-and-fruit}{The fruit of a gift}
God's demand for fruit begins with what he has given. The vineyard exists because someone planted it; the disciples bear fruit because Christ chose and appointed them; the communicants receive nourishment because Christ gives himself. Yet gift does not make response optional. Isaiah's accusation and the tenants' judgment show how seriously the giver regards the life made possible by his generosity. The first reading of the whole formulary holds these two truths together: obedience is not the purchase price of God's initiative, and grace does not make injustice harmless.

\subsection{Before the harvest there is a planting}
Chrysostom starts his exposition of the tenants with the actions of the householder. Planting, enclosure, winepress and tower show that the owner has taken care of the vineyard before demanding its produce. The work has a history prior to the tenants' arrival. They enter an already prepared place, and their task is to care for what they received. For Chrysostom these provisions recall God's care in bringing the people from Egypt and giving them law, city, temple and altar. The image records an inheritance of benefits, not an isolated contract between equally independent parties.\footnote{Chrysostom, \work{Homilies on Matthew} 68.1, NPNF1, vol. 10.}

This is why the violence is so serious. The messengers do not demand something alien to the tenancy; they ask for the fruit toward which the preparation was directed. Chrysostom identifies it as obedience evidenced by works. The tenants turn the relationship inside out. Rather than acknowledge a debt of grateful service, they punish the people who name it. Their cruelty protects a false picture of themselves as possessors answerable to nobody. The final attack upon the heir is the fullest expression of that picture.

The owner's patience then reveals both his generosity and their culpability. He sends again after injury. Finally he sends his son. Chrysostom reads the expectation that they will respect the son as what they ought to do, preserving both God's knowledge and the reality of their refusal. Knowledge of a sinful act does not make the command to receive the son meaningless. The account therefore offers no comfortable determinism in which everyone simply acts out an unavoidable role. The hearer is asked to recognize a gift and answer it.

Isaiah begins still earlier in the image: the vineyard itself fails to yield the expected grapes. Its final identification of fruit with justice supplies a concrete test of that failure. Chrysostom's distinction between vineyard and husbandmen keeps the two accusations from collapsing into each other. In Isaiah, the planting names Israel and Judah; in Matthew the immediate accusation focuses upon their rulers. Together the passages expose both a people's unjust life and the violence of stewards who resist correction. Neither layer can be repaired by the mere continuation of religious appearances.

The Entrance reaches this same dependence through a different experience. Mordecai invokes the Creator from within a threat he cannot master. The danger does not prove a failure of merit; it occasions intercession. This opening voice gives the vineyard image its proper scale. Even a richly cultivated people does not own the world, and even a powerless people remains able to turn toward its Maker. Divine sovereignty is not an attribute the tenants can appropriate for themselves. It is the ground on which the endangered can pray.

\subsection{Chosen so as to become good}
Augustine's exposition of John 15:16 asks what Christ found in those he chose. Could their prior goodness have supplied the reason for his choice? Augustine answers by following the order of the saying. If people had first turned to Christ in such a way as to establish their worthiness, their choice would have come first. Christ says instead that his choice precedes theirs. Augustine therefore understands mercy as producing the goodness it does not merely discover.\footnote{Augustine, \work{Tractates on John} 86.2--3, trans. John Gibb and James Innes, NPNF1, vol. 7.}

His next step is just as necessary. Chosen people are appointed to go and bear fruit. The priority of grace is not an exemption from goodness; it is its enabling source. Augustine recalls the branch's inability to bear apart from the vine, making continued dependence part of the very act of fruitfulness. A branch does not graduate from receiving life into producing independently. Christ is also the way in which those appointed are to go. The gift reaches beyond an initial moment to sustain the journey it begins.

This argument changes the way the Sunday can be heard. Isaiah's owner asks why good care has not yielded good fruit; the Gospel threatens unfruitful stewardship; the acclamation directs the listener to Christ's appointment. If these demands were heard without the prior gift, they could become an anxious effort to establish a claim upon God. Augustine forbids that inversion. The Christian acts because Christ has given life and a calling, not in order to make a first gift owed. The seriousness of action survives, while boasting loses its basis.

The Collect gives a corresponding form to the Church's petition. Divine generosity exceeds merit and the measure of asking. Even the prayer's confidence comes from God rather than from a calculation that the applicant has qualified. Forgiveness and unearned abundance stand together. The Collect can thus accompany the vineyard's severe accusation without allowing either evasion or hopelessness: wrongdoing is real enough to require mercy, and mercy is generous enough to be sought by the guilty.

Chrysostom's exhortation in the Philippians homily has a different accent. Near its end he urges the hearers to make a beginning and cease fighting. He is summoning people to act; Augustine is explaining why prior goodness cannot be the cause of Christ's election. Chrysostom has already grounded peace in God's astonishing beneficence toward enemies. The compatible claim is that divine generosity calls forth active response. Their different questions should remain audible, rather than being fused into one technical account of every relation between grace and human initiative.\footnote{Chrysostom, \work{Homilies on Philippians} 14, especially the exposition of 4:7 and the final exhortation to peace.}

\subsection{The fruit has a name}
The command following John 15:16 tells Augustine what the fruit is: love of one another. His argument does not stop with an attractive image of growth. He reads the next sentence as an explanation of the previous one. The branch's fruit is charity, and charity includes love of God together with rightly ordered love of neighbor. The several virtues cannot replace charity as isolated achievements. They arise within it: patience continues in a loved good, kindness aids a loved person, and peace belongs with those one truly loves.\footnote{Augustine, \work{Tractates on John} 87.1.}

That identification brings the Gospel's demand into practical contact with Isaiah's justice. The outcry of someone wronged cannot be dismissed by pointing to impressive religious activity elsewhere. Charity seeks the person's good; justice names what is owed. In this reading they expose the inadequacy of fruit understood as quantities alone. A person can increase possessions, influence or work accomplished while remaining unwilling to receive correction or serve a neighbor. The tenants' desire to acquire the inheritance is the parable's most obvious example of growth that is actually ruin.

Philippians supplies the ordinary practices in which an opposite life develops. Requests are accompanied by thanksgiving, which remembers that petitioners already live amid gifts. Attention is turned toward truth, justice, purity and what deserves praise. Teaching becomes action through the example of a life seen and heard. Chrysostom treats the transition from thought to deed as essential: evil actions have sources in the things a person entertains, while the apostolic example gives the hearer a pattern to practice. Fruit is cultivated in repeated acts rather than secured by one statement of good intention.

The peace Paul promises is also more than the reward of having managed oneself efficiently. Chrysostom relates it to God's reconciliation of those who were hostile to him through the Son. The gift supplies the measure of conduct toward other people. If God did not restrict his generosity to those who had earned his friendship, the hearer cannot make personal grievance the supreme rule of action. Thanksgiving and forgiveness both acknowledge that the self is not the original source of the goods it receives.

\subsection{A prayer for life rather than self-sufficiency}
The psalm turns demand into petition without withdrawing the demand. The people ask God to give life, and then they will invoke his name. In Augustine's closing exposition, quickening changes what the heart loves: the Creator is no longer valued merely as a supplier of earthly benefits. The living relation to God becomes itself the good sought. This gives the vineyard's fruit a deeper horizon than visible prosperity.\footnote{Augustine, \work{Exposition of Psalm} 79, §11, NPNF1, vol. 8, heading Psalm 80.}

Augustine asks what created thing could rightly be loved more than its Creator. His answer proceeds through the beauty of the world rather than contempt for it. Created beauty invites love of the one who made it. The danger comes when a gift becomes the preferred object that displaces the giver. His examples of wealth, health and honor sharpen the point: those goods are also possessed by people who do not worship God. If they become the sole reason for worship, the worshipper will judge God by the prosperity of others and lose the purpose of love.

Lamentations' Communion option holds this dependence within grief. Seeking the good Lord does not presuppose that all consolation has returned. The surrounding chapter allows brokenness, memory, waiting and hope to occupy the same human life. Fruit may therefore include a fidelity whose full outcome is not yet seen. Such waiting is not passive approval of injustice; the same chapter protests violated rights and calls for examination and return. It directs hope toward God without making present pain unreal.

The Corinthians alternative draws out a different consequence. If the many share one bread and become one body, the neighbor is not merely another independent worker attempting to deliver a separate crop. Charity bears fruit within a shared sacramental life. Chrysostom's question is whether those nourished by the same Christ actually show the same love. The sacrament reveals both the source of unity and the contradiction of behavior that divides its recipients.\footnote{Chrysostom, \work{Homilies on First Corinthians} 24.4.}

\subsection{The gift nourishes the response}
Aquinas explains the Eucharist's gift by beginning with Christ, whom the sacrament contains, and with the Passion it represents. The spiritual benefit is not inferred merely from bread as a useful illustration. Christ gives life, and his Passion is the source of the sacrament's efficacy. Nourishment then explains how that life is sustained, increased, restored and delighted. The outward species also signify unity: many grains and grapes are gathered into one.\footnote{Aquinas, \work{Summa theologiae} III, q. 79, a. 1, body. Historical English translation by the Fathers of the English Dominican Province.}

The Prayer over the Offerings and Prayer after Communion occupy distinct places within that action. The former joins obedient worship to redemption's sanctifying effect. The latter seeks the shaping of the receiver by the gift received. Neither is a promise that ritual performance purchases an unrestricted favor. In Aquinas's reply to the second objection, charity is not only bestowed as a disposition but stirred into action. The sacrament strengthens precisely the love that Augustine has identified as fruit. The demand and the nourishment meet without becoming identical acts.

This account also answers a difficulty in the language of gift. If everything begins with God, does the person disappear into passivity? Aquinas's description is the opposite: charity is moved to work. The grace received reaches embodied action, and the body is offered in the service of justice. The point is not to replace action with a claim that God will act instead. It is to understand action as the living expression of a gift whose source remains Christ.\footnote{Aquinas, III, q. 79, a. 1, replies 2--3.}

\subsection{Fruit that remains}
The harvest image could suggest only one season's results. John speaks of fruit remaining, and Augustine relates present love to desire for a fulfillment not yet possessed. His exposition of the psalm makes God himself the reward of the justified. Aquinas adds the body's future participation in incorruption and glory. Together these accounts prevent the moral reading from closing around successful management of present life.

Judgment remains part of that horizon. The tenants cannot make the inheritance theirs by excluding the Son, and fruit remains the criterion after their removal. Hope is therefore not confidence that present membership makes perseverance unnecessary. It is trust in the giver who can sustain the response he asks, joined to the willingness to receive his correction. The Church prays as a vineyard in need of life while seeking the charity that will remain.

\begin{fourSenses}
\item[Literal.] Distinct speakers and situations express dependence and responsibility: endangered Israel appeals to its Creator; Isaiah names injustice; the devastated vine asks for life; Paul teaches thankful prayer and practice; Jesus judges violent stewardship. The prayers place need and response before divine generosity.
\item[Allegorical.] The rejected Son is the cornerstone, and the disciples' fruit grows through abiding in Christ. Augustine's charity and Chrysostom's one body disclose the unity of giver, nourishment and fruitful life in Christ and his members.
\item[Moral.] Gifts are received as a call to justice, truth, thanksgiving and mutual service. Works answer Christ's initiative; they do not purchase his initial choice. Repentance receives correction where possessiveness has withheld fruit.
\item[Anagogical.] Lasting charity tends toward enjoyment of God himself and the body's future incorruption. The harvest's judgment makes perseverance serious; the giver's faithfulness makes hope possible beyond the measure of one visible season.
\end{fourSenses}
